\FigureLayoutDeclare{img/2014/new_curriculum_paper_1_liberal/q22_fig1.png}{8.1cm}{6bp 0bp 71bp 15bp}

\chapter{2014年新课标I卷（文）}

\section{单选题}

\begin{problem}
已知集合 \(M = \{x \mid -1 < x < 3\}\)，\(N = \{x \mid -2 < x < 1\}\)，则 \(M \cap N =\)（\quad）

\choices
    {\((-2, 1)\)}
    {\((-1, 1)\)}
    {\((1,3)\)}
    {\((-2,3)\)}
\end{problem}

\begin{answer}
B
\end{answer}

\begin{solution}
交集中的元素必须同时属于 \(M\) 和 \(N\)，所以必须满足 \(-1 < x < 3\) 且 \(-2 < x < 1\). 合并两个不等式得 \(-1 < x < 1\)，因此 \(M \cap N = (-1,1)\)，故选 B.
\end{solution}

\begin{problem}
若 \(\tan \alpha > 0\)，则（\quad）

\choices
    {\(\sin 2\alpha > 0\)}
    {\(\cos \alpha > 0\)}
    {\(\sin \alpha > 0\)}
    {\(\cos 2\alpha > 0\)}
\end{problem}

\begin{answer}
A
\end{answer}

\begin{solution}
由 \(\tan\alpha=\dfrac{\sin\alpha}{\cos\alpha}>0\)，知 \(\sin\alpha\) 与 \(\cos\alpha\) 同号，且二者均不为零. 因此 \(\sin 2\alpha=2\sin\alpha\cos\alpha>0\). \(\tan\alpha>0\) 时 \(\alpha\) 也可能在第三象限，故不能确定 \(\sin\alpha\) 或 \(\cos\alpha\) 的单独符号；同样 \(\cos 2\alpha\) 的符号也不能确定. 故选 A.
\end{solution}

\begin{problem}
设 \(z = \frac{1}{1 + \i} + \i\)，则 \(|z| =\)（\quad）

\choices
    {\(\frac{1}{2}\)}
    {\(\frac{\sqrt{2}}{2}\)}
    {\(\frac{\sqrt{3}}{2}\)}
    {\(2\)}
\end{problem}

\begin{answer}
B
\end{answer}

\begin{solution}
分母有理化得 \(\dfrac{1}{1+\i}=\dfrac{1-\i}{(1+\i)(1-\i)}=\dfrac{1-\i}{2}\)，所以 \(z=\dfrac{1-\i}{2}+\i=\dfrac{1+\i}{2}\). 因此 \(|z|=\sqrt{\left(\dfrac{1}{2}\right)^2+\left(\dfrac{1}{2}\right)^2}=\dfrac{\sqrt{2}}{2}\)，故选 B.
\end{solution}

\begin{problem}
已知双曲线 \(\frac{x^2}{a^2} - \frac{y^2}{3} = 1\)（\(a > 0\)）的离心率为 \(2\)，则 \(a =\)（\quad）

\choices
    {\(2\)}
    {\(\frac{\sqrt{6}}{2}\)}
    {\( \frac{\sqrt{5}}{2} \)}
    {\(1\)}
\end{problem}

\begin{answer}
D
\end{answer}

\begin{solution}
双曲线的标准式中 \(a^2\) 为横轴半轴平方，\(b^2=3\)，所以 \(c^2=a^2+b^2=a^2+3\). 由离心率 \(e=\dfrac{c}{a}=2\)，得 \(a^2+3=4a^2\)，从而 \(a^2=1\). 又 \(a>0\)，所以 \(a=1\)，故选 D.
\end{solution}

\begin{problem}
设函数 \( f(x)\)，\(g(x)\) 的定义域都为 \( \R \)，且 \( f(x) \) 是奇函数，\( g(x) \) 是偶函数，则下列结论正确的是（\quad）

\choices
    {\(f(x)|g(x)|\) 是奇函数}
    {\( |f(x)|g(x) \) 是奇函数}
    {\(f(x)g(x)\) 是偶函数}
    {\( |f(x)g(x)| \) 是奇函数}
\end{problem}

\begin{answer}
    A
\end{answer}

\begin{solution}
因为 \(f\) 是奇函数，\(f(-x)=-f(x)\)；因为 \(g\) 是偶函数，\(g(-x)=g(x)\)，从而 \(|g(-x)|=|g(x)|\). 于是 \([f(x)|g(x)|]_{x\mapsto -x}=f(-x)|g(-x)|=-f(x)|g(x)|\)，所以 \(f(x)|g(x)|\) 是奇函数. 另一方面，\(|f(x)|g(x)\) 以及 \(|f(x)g(x)|\) 都是偶函数，而 \(f(x)g(x)\) 是奇函数. 故正确结论为 A.
\end{solution}

\begin{problem}
设 \(D\)，\(E\)，\(F\) 分别为 \(\triangle ABC\) 的三边 \(BC\)，\(CA\)，\(AB\) 的中点，则 \(\overrightarrow{EB} + \overrightarrow{FC} =\)（\quad）

\choices
    {\(\overrightarrow{BC}\)}
    {\(\frac{1}{2}\overrightarrow{AD}\)}
    {\(\overrightarrow{AD}\)}
    {\(\frac{1}{2}\overrightarrow{BC}\)}
\end{problem}

\begin{answer}
C
\end{answer}

\begin{solution}
设 \(A\)，\(B\)，\(C\) 的位置向量分别为 \(\bs{a}\)，\(\bs{b}\)，\(\bs{c}\). 由中点条件，\(E\) 和 \(F\) 的位置向量分别为 \(\dfrac{\bs{a}+\bs{c}}{2}\) 和 \(\dfrac{\bs{a}+\bs{b}}{2}\)，所以 \(\overrightarrow{EB}+\overrightarrow{FC}=\left(\bs{b}-\dfrac{\bs{a}+\bs{c}}{2}\right)+\left(\bs{c}-\dfrac{\bs{a}+\bs{b}}{2}\right)=\dfrac{\bs{b}+\bs{c}}{2}-\bs{a}\). 又 \(D\) 是 \(BC\) 的中点，故 \(\overrightarrow{AD}=\dfrac{\bs{b}+\bs{c}}{2}-\bs{a}\)，从而 \(\overrightarrow{EB}+\overrightarrow{FC}=\overrightarrow{AD}\)，故选 C.
\end{solution}

\begin{problem}
在下列函数中：

\begin{circlelist}
\item \(y = \cos |2x|\)；
\item \(y = |\cos x|\)；
\item \(y = \cos \left(2x + \frac{\pi}{6}\right)\)；
\item \(y = \tan \left(2x - \frac{\pi}{4}\right)\).
\end{circlelist}

最小正周期为 \(\pi\) 的所有函数为（\quad）

\choices
    {\circled{2}\circled{4}}
    {\circled{1}\circled{3}\circled{4}}
    {\circled{1}\circled{2}\circled{3}}
    {\circled{1}\circled{3}}
\end{problem}

\begin{answer}
C
\end{answer}

\begin{solution}
因为余弦函数是偶函数，\(\cos|2x|=\cos 2x\)，其最小正周期为 \(\pi\). 函数 \(|\cos x|\) 的最小正周期为 \(\pi\)，函数 \(\cos\left(2x+\dfrac{\pi}{6}\right)\) 的最小正周期为 \(\dfrac{2\pi}{2}=\pi\). 正切函数的最小正周期为 \(\pi\)，所以函数 \(\tan\left(2x-\dfrac{\pi}{4}\right)\) 的最小正周期为 \(\dfrac{\pi}{2}\). 因此满足条件的是 \(\circled{1}\)、\(\circled{2}\)、\(\circled{3}\)，故选 C.
\end{solution}

\begin{problem}
如图，网格纸的各小格都是正方形，粗实线画出的是一个几何体的三视图，则这个几何体是（\quad）

\bitmapfigure[max width=0.52\linewidth]{img/2014/new_curriculum_paper_1_liberal/q8_fig1.png}

\choices
    {三棱锥}
    {三棱柱}
    {四棱锥}
    {四棱柱}
\end{problem}

\begin{answer}
B
\end{answer}

\begin{solution}
图中的一个视图是三角形，另两个视图是矩形，且矩形表示三角形截面沿同一方向平移形成的侧面. 因此该几何体有两个全等且平行的三角形底面，其他三个面为矩形，正是三棱柱的三视图，而不是只有一个顶点的三棱锥或具有四边形底面的棱柱、棱锥. 故选 B.
\end{solution}

\begin{problem}
执行如图的程序框图，若输入的 \(a\)，\(b\)，\(k\) 分别为 \(1\)，\(2\)，\(3\)，则输出的 \(M =\)（\quad）

\bitmapfigure[max width=0.50\linewidth]{img/2014/new_curriculum_paper_1_liberal/q9_fig1.png}

\choices
    {\(\frac{20}{3}\)}
    {\(\frac{7}{2}\)}
    {\(\frac{16}{5}\)}
    {\(\frac{15}{8}\)}
\end{problem}

\begin{answer}
D
\end{answer}

\begin{solution}
初始时 \(n=1\)，\(a=1\)，\(b=2\). 第一次循环中 \(M=1+\dfrac{1}{2}=\dfrac{3}{2}\)，随后更新为 \(a=2\)，\(b=\dfrac{3}{2}\)，\(n=2\). 第二次循环中 \(M=2+\dfrac{1}{3/2}=\dfrac{8}{3}\)，随后更新为 \(a=\dfrac{3}{2}\)，\(b=\dfrac{8}{3}\)，\(n=3\). 第三次循环中 \(M=\dfrac{3}{2}+\dfrac{1}{8/3}=\dfrac{15}{8}\)，更新后 \(n=4\)，不再满足 \(n\leqslant k\). 所以输出 \(M=\dfrac{15}{8}\)，故选 D.
\end{solution}

\begin{problem}
已知抛物线 \(C: y^{2} = x\) 的焦点为 \(F\)， \(A(x_{0}, y_{0})\) 是 \(C\) 上一点，\(|AF| = \frac{5}{4} x_{0}\)，则 \(x_{0} =\)（\quad）

\choices
    {\(4\)}
    {\(2\)}
    {\(1\)}
    {\(8\)}
\end{problem}

\begin{answer}
C
\end{answer}

\begin{solution}
抛物线 \(y^2=x\) 可写成 \(y^2=4px\)，所以 \(p=\dfrac14\)，焦点为 \(F\left(\dfrac14,0\right)\). 又因 \(A\) 在抛物线上，\(y_0^2=x_0\)，且 \(x_0\geqslant0\). 于是
\(\lvert AF\rvert^2=\left(x_0-\dfrac14\right)^2+y_0^2=\left(x_0+\dfrac14\right)^2\)，从而 \(\lvert AF\rvert=x_0+\dfrac14\). 代入条件得 \(x_0+\dfrac14=\dfrac54x_0\)，解得 \(x_0=1\)，故选 C.
\end{solution}

\begin{problem}
设 \(x\)，\(y\) 满足约束条件 \(\begin{cases}
x + y \geqslant a,\\
x - y \leqslant -1,
\end{cases}\) 且 \(z = x + ay\) 的最小值为 \(7\)，则 \(a =\)（\quad）

\choices
    {\(-5\)}
    {\(3\)}
    {\(-5\)或3}
    {\(5\) 或 \(-3\)}
\end{problem}

\begin{answer}
B
\end{answer}

\begin{solution}
令 \(u=x+y\)，\(v=x-y\)，则约束变为 \(u\geqslant a\)，\(v\leqslant-1\)，且 \(x=\dfrac{u+v}{2}\)，\(y=\dfrac{u-v}{2}\). 因此 \(z=\dfrac{(a+1)u+(1-a)v}{2}\). 要使这个线性函数在上述无界区域有有限最小值，必须有 \(a+1\geqslant0\) 且 \(1-a\leqslant0\)，所以 \(a\geqslant1\). 当 \(a=1\) 时，\(z=u\) 的最小值为 \(1\)，不符合题意；故 \(a>1\). 此时最小值在 \(u=a\)，\(v=-1\) 处取得，故 \(7=\dfrac{(a+1)a+(1-a)(-1)}{2}=\dfrac{a^2+2a-1}{2}\). 于是 \(a^2+2a-15=0\)，解得 \(a=3\) 或 \(a=-5\)，结合 \(a>1\) 得 \(a=3\)，故选 B.
\end{solution}

\begin{problem}
已知函数 \( f(x) = ax^3 - 3x^2 + 1 \)，若 \( f(x) \) 存在唯一的零点 \( x_0 \)，且 \( x_0 > 0 \)，则 \( a \) 的取值范围为（\quad）

\choices
    {\((- \infty, -2)\)}
    {\((1, +\infty)\)}
    {\((2, +\infty)\)}
    {\((-\infty, -1)\)}
\end{problem}

\begin{answer}
A
\end{answer}

\begin{solution}
有 \(f(0)=1\)，且 \(f'(x)=3x(ax-2)\). 当 \(a=0\) 时，\(f(x)=-3x^2+1\) 有两个零点，不符合条件.

当 \(a>0\) 时，\(f(x)\) 在 \((-\infty,0)\) 的两端符号分别为负、正，必有负零点. 对 \(x\geqslant0\)，函数在 \(x=\dfrac{2}{a}\) 处取得最小值 \(f\left(\dfrac{2}{a}\right)=1-\dfrac{4}{a^2}\). 若 \(0<a<2\)，该最小值为负，因而有两个正零点；若 \(a=2\)，有一个正零点且另有负零点；若 \(a>2\)，没有正零点. 故 \(a>0\) 均不符合唯一零点且零点为正的条件.

当 \(a<0\) 时，\(f'(x)<0\) 对一切 \(x>0\) 成立，且 \(f(0)=1>0\)，\(f(x)\) 在 \((0,+\infty)\) 上递减并趋于负无穷，所以恰有一个正零点. 负半轴上的极小值在 \(x=\dfrac{2}{a}<0\) 处取得，且值为 \(1-\dfrac{4}{a^2}\). 当 \(-2<a<0\) 时该值小于零，产生两个负零点；当 \(a=-2\) 时产生一个负的重零点；只有当 \(a<-2\) 时极小值大于零，负半轴上没有零点. 因此所求范围为 \((-\infty,-2)\)，故选 A.
\end{solution}

\section{填空题}

\begin{problem}
将 \(2\) 本不同的数学书和 \(1\) 本语文书在书架上随机排成一行，则 \(2\) 本数学书相邻的概率为\fillinblank{}.
\end{problem}

\begin{answer}
\(\frac{2}{3}\)
\end{answer}

\begin{solution}
把两本不同的数学书和一本语文书看作三个不同对象，共有 \(3!=6\) 种排法. 两本数学书相邻时，把它们看成一个整体，整体与语文书有 \(2!\) 种排法，数学书在整体内还有 \(2!\) 种顺序，所以有 \(2!\times2!=4\) 种. 故所求概率为 \(\dfrac{4}{6}=\dfrac{2}{3}\).
\end{solution}

\begin{problem}
甲，乙，丙三位同学被问到是否去过 \(A\)，\(B\)，\(C\) 三个城市时，甲说：我去过的城市比乙多，但没去过 \(B\) 城市；乙说：我没去过 \(C\) 城市；丙说：我们三人去过同一个城市. 由此可判断乙去过的城市为\fillinblank{}.
\end{problem}

\begin{answer}
A
\end{answer}

\begin{solution}
甲没有去过 \(B\)，乙没有去过 \(C\)，而三人去过同一个城市. 这个共同去过的城市既不能是 \(B\)，也不能是 \(C\)，所以只能是 \(A\)，从而乙去过 \(A\). 若乙还去过 \(B\)，则乙至少去过两个城市；甲只能在 \(A\)，\(C\) 中选择，至多去过两个城市，不可能比乙多. 因此乙没有去过 \(B\)，又没有去过 \(C\)，所以乙去过的城市为 A.
\end{solution}

\begin{problem}
设函数 \( f(x) = \begin{cases}
\e^{x - 1}, & x < 1,\\
x^{\frac{1}{3}}, & x \geqslant 1,
\end{cases} \) 则使得 \( f(x) \leqslant 2 \) 成立的 \( x \) 的取值范围是\fillinblank{}.
\end{problem}

\begin{answer}
\(x \leqslant 8\)
\end{answer}

\begin{solution}
当 \(x<1\) 时，\(x-1<0\)，所以 \(f(x)=\e^{x-1}<1\leqslant2\)，所有 \(x<1\) 都满足条件. 当 \(x\geqslant1\) 时，\(x^{1/3}\leqslant2\) 等价于 \(x\leqslant8\)，所以这一段的解集为 \([1,8]\). 合并两段得 \((-\infty,1)\cup[1,8]=(-\infty,8]\)，即 \(x\leqslant8\).
\end{solution}

\begin{problem}
如图，为测量山高 \(MN\)，选择 \(A\) 和另一座山的山顶 \(C\) 为测量观测点. 从 \(A\) 点测得 \(M\) 点的仰角 \(\angle MAN = 60^\circ\)，\(C\) 点的仰角 \(\angle CAB = 45^\circ\) 以及 \(\angle MAC = 75^\circ\)；从 \(C\) 点测得 \(\angle MCA = 60^\circ\). 已知山高 \(BC = 100 \mathrm{~m}\)，则山高 \(MN =\)\fillinblank{}\(\mathrm{m}\).

\bitmapfigure[max width=0.38\linewidth]{img/2014/new_curriculum_paper_1_liberal/q16_fig1.png}
\end{problem}

\begin{answer}
\(150\)
\end{answer}

\begin{solution}
在直角三角形 \(ABC\) 中，\(\angle CAB=45^\circ\)，\(BC=100\mathrm{~m}\)，所以 \(\angle ABC=90^\circ\)，且 \(AC=\dfrac{BC}{\sin45^\circ}=100\sqrt{2}\mathrm{~m}\). 在三角形 \(AMC\) 中，\(\angle MAC=75^\circ\)，\(\angle MCA=60^\circ\)，故 \(\angle AMC=45^\circ\). 由正弦定理，\(\dfrac{AM}{\sin60^\circ}=\dfrac{AC}{\sin45^\circ}\)，得 \(AM=100\sqrt{3}\mathrm{~m}\). 又在直角三角形 \(AMN\) 中，\(MN=AM\sin60^\circ=100\sqrt{3}\times\dfrac{\sqrt{3}}{2}=150\mathrm{~m}\)，所以填 \(150\).
\end{solution}

\section{解答题}

\begin{problem}
已知 \(\{a_{n}\}\) 是递增的等差数列，\(a_{2}\)，\(a_{4}\) 是方程 \(x^{2} - 5x + 6 = 0\) 的根.

\begin{enumerate}
\item 求 \(\{a_{n}\}\) 的通项公式；
\item 求数列 \(\left\{\frac{a_{n}}{2^{n}}\right\}\) 的前 \(n\) 项和.
\end{enumerate}
\end{problem}

\begin{solution}
\begin{enumerate}
\item 方程 \(x^2-5x+6=0\) 的两个根为 \(2\)，\(3\). 数列递增，所以 \(a_2=2\)，\(a_4=3\). 设公差为 \(d\)，则 \(a_4-a_2=2d=1\)，得 \(d=\dfrac12\). 于是 \(a_n=a_2+(n-2)d=2+\dfrac{n-2}{2}=\dfrac n2+1\).
\item 由上式，\(\dfrac{a_k}{2^k}=\dfrac{k+2}{2^{k+1}}\)，所以
\[S_n=\sum_{k=1}^{n}\frac{k+2}{2^{k+1}}=\frac12\sum_{k=1}^{n}\frac{k}{2^k}+\sum_{k=1}^{n}\frac1{2^k}\]
令 \(T_n=\sum_{k=1}^{n}\dfrac{k}{2^k}\)，则
\[T_n=\frac12+\frac{2}{2^2}+\cdots+\frac{n}{2^n}\]
两边同乘 \(\dfrac12\)，得
\[\frac12T_n=\frac1{2^2}+\frac{2}{2^3}+\cdots+\frac{n}{2^{n+1}}\]
两式相减得
\[\frac12T_n=\frac12+\sum_{k=2}^{n}\frac1{2^k}-\frac{n}{2^{n+1}}=1-\frac1{2^n}-\frac{n}{2^{n+1}}\]
所以 \(T_n=2-\dfrac{n+2}{2^n}\). 另外，由等比数列求和公式，\(\sum_{k=1}^{n}\dfrac1{2^k}=1-\dfrac1{2^n}\). 因此
\[S_n=\frac12\left(2-\frac{n+2}{2^n}\right)+1-\frac1{2^n}=2-\frac{n+4}{2^{n+1}}\]
\end{enumerate}
\end{solution}

\begin{problem}
从某企业生产的某种产品中抽取 \(100\) 件，测量这些产品的一项质量指标值，由测量结果得如下频数分布表：

\begin{center}
\begin{tabular}{c|ccccc}
\hline
质量指标值分组 & \([75,85)\) & \([85,95)\) & \([95,105)\) & \([105,115)\) & \([115,125)\) \\
\hline
频数 & \(6\) & \(26\) & \(38\) & \(22\) & \(8\) \\
\hline
\end{tabular}
\end{center}

\begin{enumerate}
\item 作出这些数据的频率分布直方图；
\item 估计这种产品质量指标值的平均数及方差（同一组中的数据用该组区间的中点值作代表）；
\item 根据以上抽样调查数据，能否认为该企业生产的这种产品符合“质量指标值不低于 \(95\) 的产品至少要占全部产品的 \(80\%\)”的规定？
\end{enumerate}

\bitmapfigure[max width=0.38\linewidth]{img/2014/new_curriculum_paper_1_liberal/q18_fig1.png}
\FloatBarrier
\end{problem}

\begin{solution}
\begin{enumerate}
\item 各组组距均为 \(10\)，所以频率分布直方图中各矩形的高为“频率除以组距”. 五组频率分别为 \(0.06\)，\(0.26\)，\(0.38\)，\(0.22\)，\(0.08\)，故各组柱高依次为 \(0.006\)，\(0.026\)，\(0.038\)，\(0.022\)，\(0.008\).
\item 各组中点依次为 \(80\)，\(90\)，\(100\)，\(110\)，\(120\). 平均数的估计值为 \(80\times0.06+90\times0.26+100\times0.38+110\times0.22+120\times0.08=100\). 以 \(100\) 为平均数，方差的估计值为 \(20^2\times0.06+10^2\times0.26+0^2\times0.38+10^2\times0.22+20^2\times0.08=104\).
\item 质量指标值不低于 \(95\) 的样本占比为 \(\dfrac{38+22+8}{100}=0.68\)，小于规定的 \(0.80\). 因此根据抽样调查数据，不能认为该企业生产的产品符合规定.
\end{enumerate}
\end{solution}

\begin{problem}
如图，三棱柱 \(ABC - A_{1}B_{1}C_{1}\) 中，侧面 \(BB_{1}C_{1}C\) 为菱形，\(B_{1}C\) 的中点为 \(O\)，且 \(AO \perp\) 平面 \(BB_{1}C_{1}C\).

\begin{enumerate}
\item 证明：\(B_{1}C \perp AB\)；
\item 若 \(AC \perp AB_{1}\)，\(\angle CBB_{1} = 60^{\circ}\)，\(BC = 1\)，求三棱柱 \(ABC - A_{1}B_{1}C_{1}\) 的高.
\end{enumerate}

\bitmapfigure[float=inline,max width=0.38\linewidth]{img/2014/new_curriculum_paper_1_liberal/q19_fig1.png}
\FloatBarrier
\end{problem}

\begin{solution}
\begin{enumerate}
\item 菱形的两条对角线互相垂直且互相平分. 由 \(O\) 是 \(B_1C\) 的中点，且 \(O\) 也是另一条对角线 \(BC_1\) 的中点，得 \(B_1C\perp BC_1\). 又因 \(AO\perp\) 平面 \(BB_1C_1C\)，所以 \(AO\perp B_1C\). 由于 \(O\) 在 \(BC_1\) 上，直线 \(BO\) 与 \(BC_1\) 重合，因此直线 \(AO\) 与 \(BO\) 相交于 \(O\)，且二者都垂直于 \(B_1C\). 故 \(B_1C\perp\) 平面 \(ABO\)，从而 \(B_1C\perp AB\).
\item 取侧面所在平面为 \(xy\) 平面，以两条互相垂直的菱形对角线为坐标轴，并令 \(O\) 为原点. 设 \(B_1=(p,0,0)\)，\(C=(-p,0,0)\)，\(B=(0,-q,0)\)，则 \(p\)，\(q>0\)，且 \(p^2+q^2=BC^2=1\). 由 \(\angle CBB_1=60^\circ\)，有
\[\frac{(-p,q,0)\cdot(p,q,0)}{1\times1}=q^2-p^2=\cos60^\circ=\frac12\]
故 \(p^2=\dfrac14\)，\(q^2=\dfrac34\). 由 \(AO\perp\) 该平面，设 \(A=(0,0,h)\)，其中 \(h>0\). 条件 \(AC\perp AB_1\) 给出 \((-p,0,-h)\cdot(p,0,-h)=h^2-p^2=0\)，所以 \(h=\dfrac12\).

三棱柱的另一底面由平移向量 \(\overrightarrow{BB_1}=(p,q,0)\) 得到. 平面 \(ABC\) 内的两个方向向量可取 \(\overrightarrow{BC}=(-p,q,0)\)，\(\overrightarrow{BA}=(0,q,h)\)，其一个法向量为 \(\bs{n}=(qh,ph,-pq)\). 因此两底面间的距离，即三棱柱的高为
\[H=\frac{\left|\overrightarrow{BB_1}\cdot\bs{n}\right|}{|\bs{n}|}=\frac{2pqh}{\sqrt{h^2(p^2+q^2)+p^2q^2}}=\frac{\sqrt3}{\sqrt7}=\frac{\sqrt{21}}{7}\]
\end{enumerate}
\end{solution}

\begin{problem}
已知点 \(P(2,2)\)，圆 \(C: x^{2} + y^{2} - 8y = 0\)，过点 \(P\) 的动直线 \(l\) 与圆 \(C\) 交于 \(A\)，\(B\) 两点. 线段 \(AB\) 的中点为 \(M\)，\(O\) 为坐标原点.

\begin{enumerate}
\item 求 \(M\) 的轨迹方程；
\item 当 \(|OP| = |OM|\) 时，求 \(l\) 的方程及 \(\triangle POM\) 的面积.
\end{enumerate}
\end{problem}

\begin{solution}
\begin{enumerate}
\item 圆 \(C\) 的方程为 \(x^2+(y-4)^2=16\)，圆心记为 \(O_1=(0,4)\). 设 \(M=(x,y)\)，则圆心与弦中点的连线垂直于弦，故 \(\overrightarrow{O_1M}\perp\overrightarrow{PM}\). 于是 \((x,y-4)\cdot(x-2,y-2)=0\)，即 \(x(x-2)+(y-4)(y-2)=0\)，整理得 \((x-1)^2+(y-3)^2=2\). 反之，该圆的圆心为 \(N=(1,3)\)，半径为 \(\sqrt{2}\)，且 \(\lvert O_1N\rvert+\sqrt{2}=2\sqrt{2}<4\)，所以圆上的每个点都在原圆内部. 当 \(M\ne P\) 时，等式给出 \(O_1M\perp PM\)，故直线 \(PM\) 与原圆的交点以 \(M\) 为中点；当 \(M=P\) 时，过 \(P\) 作垂直于 \(O_1P\) 的直线，该直线与原圆交于两个点且以 \(P\) 为弦中点. 因此 \(M\) 的轨迹方程为 \((x-1)^2+(y-3)^2=2\).
\item 记轨迹圆的圆心为 \(N=(1,3)\). 由于题目要求 \(\triangle POM\)，取非退化情形 \(M\ne P\)；若允许退化情形，则 \(M=P\) 时的直线为 \(y=x\)，所得三点不构成题目所称的三角形. 由 \(|OP|=|OM|\)，点 \(O\) 在 \(PM\) 的垂直平分线上；又因 \(P\)，\(M\) 都在以 \(N\) 为圆心的轨迹圆上，点 \(N\) 也在该垂直平分线上. 因为 \(P\ne M\)，这条垂直平分线就是直线 \(ON\)，所以 \(ON\perp PM\). 直线 \(ON\) 的斜率为 \(3\)，故直线 \(l=PM\) 的斜率为 \(-\dfrac13\)，且过 \(P(2,2)\)，从而 \(l\) 的方程为 \(x+3y-8=0\).

\(|OP|=2\sqrt2\)，点 \(O\) 到直线 \(l\) 的距离为 \(\dfrac8{\sqrt{10}}=\dfrac{4\sqrt{10}}5\)，点 \(N\) 到直线 \(l\) 的距离为 \(\dfrac2{\sqrt{10}}=\dfrac{\sqrt{10}}5\). 因此弦 \(PM\) 的长为 \(2\sqrt{2-\left(\dfrac{\sqrt{10}}5\right)^2}=\dfrac{4\sqrt{10}}5\). 所以 \(\triangle POM\) 的面积为 \(\dfrac12\times\dfrac{4\sqrt{10}}5\times\dfrac{4\sqrt{10}}5=\dfrac{16}{5}\).
\end{enumerate}
\end{solution}

\begin{problem}
设函数 \( f(x) = a \ln x + \frac{1 - a}{2} x^2 - bx\)（\(a \neq 1\)），曲线 \( y = f(x) \) 在点 \((1,f(1))\) 处的切线斜率为 \(0\).

\begin{enumerate}
\item 求 \(b\)；
\item 若存在 \(x_0 \geqslant 1\) 使得 \(f(x_0) < \frac{a}{a - 1}\)，求 \(a\) 的取值范围.
\end{enumerate}
\end{problem}

\begin{solution}
\begin{enumerate}
\item 函数定义域为 \(x>0\)，且 \(f'(x)=\dfrac{a}{x}+(1-a)x-b\). 切线斜率为 \(0\) 给出 \(f'(1)=a+1-a-b=0\)，所以 \(b=1\).
\item 代入 \(b=1\)，得 \(f(x)=a\ln x+\dfrac{1-a}{2}x^2-x\)，并且
\[f'(x)=\frac{a+(1-a)x^2-x}{x}=\frac{(x-1)((1-a)x-a)}{x}\]
当 \(a>1\) 时，对一切 \(x>1\) 有 \((1-a)x-a<0\)，故 \(f\) 在 \([1,+\infty)\) 上递减且趋于负无穷，必存在满足条件的 \(x_0\).

当 \(a<1\) 时，\(f(x)\) 在 \(x\to+\infty\) 时趋于正无穷. 若 \(a\leqslant\dfrac12\)，则 \((1-a)x-a\geqslant1-2a\geqslant0\)（\(x\geqslant1\)），所以 \(f\) 在 \([1,+\infty)\) 上递增，最小值为 \(f(1)=-\dfrac{a+1}{2}\). 此时存在所求 \(x_0\) 的条件为
\[f(1)-\frac{a}{a-1}=\frac{1-2a-a^2}{2(a-1)}<0\]
由于 \(a-1<0\)，这等价于 \(a^2+2a-1<0\)，即 \(-\sqrt2-1<a<\sqrt2-1\).

若 \(\dfrac12<a<1\)，则 \(f\) 在 \(x_* =\dfrac{a}{1-a}>1\) 处取得最小值. 直接代入得
\[f(x_*)-\frac{a}{a-1}=a\left(\ln\frac{a}{1-a}+\frac{a}{2(1-a)}\right)>0\]
所以不存在满足条件的 \(x_0\). 结合 \(a\neq1\)，所求范围为 \((-\sqrt2-1,\sqrt2-1)\cup(1,+\infty)\).
\end{enumerate}
\end{solution}

\section{选考题：解答题}

\begin{problem}
如图，四边形 \(ABCD\) 是 \(\odot O\) 的内接四边形，\(AB\) 的延长线与 \(DC\) 的延长线交于点 \(E\)，且 \(CB = CE\).

\begin{enumerate}
\item 证明：\(\angle D = \angle E\)；
\item 设 \(AD\) 不是 \(\odot O\) 的直径，\(AD\) 的中点为 \(M\)，且 \(MB = MC\)，证明 \(\triangle ADE\) 为等边三角形.
\end{enumerate}

\bitmapfigure[float=inline,max width=0.36\linewidth]{img/2014/new_curriculum_paper_1_liberal/q22_fig1.png}
\FloatBarrier
\end{problem}

\begin{solution}
\begin{enumerate}
\item 设 \(\angle ADC=\delta\). 四边形 \(ABCD\) 内接，所以对角互补，\(\angle ABC+\angle ADC=180^\circ\). 又因 \(A\)，\(B\)，\(E\) 共线，\(\angle CBE=180^\circ-\angle ABC=\delta\). 在等腰三角形 \(CBE\) 中，\(CB=CE\)，故 \(\angle BEC=\angle CBE=\delta\). 而 \(D\)，\(C\)，\(E\) 共线，题图中的 \(\angle E\) 即 \(\angle BEC\)，从而 \(\angle D=\angle E\).
\item 设 \(O\) 为圆心. 由于 \(OB=OC\) 且 \(MB=MC\)，点 \(O\)，\(M\) 都在 \(BC\) 的垂直平分线上. 若 \(O=M\)，则圆心在弦 \(AD\) 上，\(AD\) 是直径，与条件矛盾，所以 \(O\ne M\)，从而 \(OM\perp BC\). 又因 \(M\) 是弦 \(AD\) 的中点，\(OM\perp AD\)，故 \(AD\parallel BC\).

由第一问设 \(\angle D=\angle E=\delta\). 在等腰三角形 \(BCE\) 中，\(\angle CBE=\angle BEC=\delta\)，所以 \(\angle BCE=180^\circ-2\delta\). 因 \(AD\parallel BC\)，且 \(D\)，\(C\)，\(E\) 共线，有 \(\angle ADE=\angle BCE=180^\circ-2\delta\). 另一方面，\(D\)，\(C\)，\(E\) 共线又给出 \(\angle ADE=\angle ADC=\delta\)，所以 \(\delta=180^\circ-2\delta\)，即 \(\delta=60^\circ\). 再由 \(AD\parallel BC\) 和 \(A\)，\(B\)，\(E\) 共线，得 \(\angle DAE=\angle CBE=60^\circ\). 于是三角形 \(ADE\) 的两个内角均为 \(60^\circ\)，第三个内角也为 \(60^\circ\)，所以 \(\triangle ADE\) 为等边三角形.
\end{enumerate}
\end{solution}

\begin{problem}
已知曲线 \(C: \frac{x^2}{4} + \frac{y^2}{9} = 1\)，直线 \(l: \begin{cases}
x = 2 + t,\\
y = 2 - 2t,
\end{cases}\) （\(t\) 为参数）.

\begin{enumerate}
\item 写出曲线 \(C\) 的参数方程，直线 \(l\) 的普通方程；
\item 过曲线 \(C\) 上任一点 \(P\) 作与 \(l\) 夹角为 \(30^{\circ}\) 的直线，交 \(l\) 于点 \(A\)，求 \(|PA|\) 的最大值与最小值.
\end{enumerate}
\end{problem}

\begin{solution}
\begin{enumerate}
\item 椭圆的参数方程可取 \(x=2\cos\theta\)，\(y=3\sin\theta\)，其中 \(\theta\) 为参数. 由直线参数方程得 \(t=x-2\)，代入 \(y=2-2t\) 得 \(y=6-2x\)，所以直线的普通方程为 \(2x+y-6=0\).
\item 设 \(P=(2\cos\theta,3\sin\theta)\). 点 \(P\) 到直线 \(l\) 的距离为
\[d=\frac{|4\cos\theta+3\sin\theta-6|}{\sqrt5}\]
因为 \(4\cos\theta+3\sin\theta\leqslant\sqrt{4^2+3^2}=5<6\)，所以 \(d=\dfrac{6-4\cos\theta-3\sin\theta}{\sqrt5}\). 直线 \(PA\) 与 \(l\) 的夹角为 \(30^\circ\)，故 \(d=|PA|\sin30^\circ=\dfrac{|PA|}{2}\)，从而 \(|PA|=\dfrac{12-8\cos\theta-6\sin\theta}{\sqrt5}\). 又 \(-10\leqslant8\cos\theta+6\sin\theta\leqslant10\)，所以 \(\dfrac{2}{\sqrt5}\leqslant|PA|\leqslant\dfrac{22}{\sqrt5}\). 因此最大值为 \(\dfrac{22\sqrt5}{5}\)，最小值为 \(\dfrac{2\sqrt5}{5}\).
\end{enumerate}
\end{solution}

\begin{problem}
若 \(a > 0\)，\(b > 0\)，且 \(\frac{1}{a} + \frac{1}{b} = \sqrt{ab}\).

\begin{enumerate}
\item 求 \(a^{3} + b^{3}\) 的最小值；
\item 是否存在 \(a\)，\(b\)，使得 \(2a + 3b = 6\)？并说明理由.
\end{enumerate}
\end{problem}

\begin{solution}
\begin{enumerate}
\item 令 \(p=\sqrt{ab}>0\). 原条件等价于 \(a+b=(ab)^{3/2}=p^3\). 由基本不等式 \(a+b\geqslant2\sqrt{ab}=2p\)，得 \(p^2\geqslant2\)，即 \(p\geqslant\sqrt2\). 又
\[a^3+b^3=(a+b)^3-3ab(a+b)=p^9-3p^5\]
在 \(p\geqslant\sqrt2\) 时，函数 \(p^9-3p^5\) 的导数为 \(3p^4(3p^4-5)>0\)，所以最小值在 \(p=\sqrt2\) 时取得. 此时 \(a+b=2\sqrt2\)，\(ab=2\)，由等号条件得 \(a=b=\sqrt2\)，故 \(a^3+b^3=4\sqrt2\).
\item 若另有 \(2a+3b=6\)，则 \(b=\dfrac{6-2a}{3}\)，且 \(0<a<3\). 于是
\[ab=2a-\frac23a^2=\frac32-\frac23\left(a-\frac32\right)^2\leqslant\frac32<2\]
但由原条件已经推出 \(ab=p^2\geqslant2\)，矛盾. 因此不存在满足条件的正数 \(a\)，\(b\).
\end{enumerate}
\end{solution}
